%========================================================================================
%	摘要
%========================================================================================
\begin{center}
{\zihao{4}\bfseries 摘\quad 要\par}
\end{center}
\begin{quotation}

\textbf{研究背景与目的：}在城市化持续推进的背景下，噪声污染已成为影响居民生活质量与社会治理成效的重要问题。本报告以浙江省杭州、宁波、温州三地常住居民为调查对象，整合省级年度噪声监测数据、政务投诉文本、问卷与访谈资料，系统分析区域噪声态势、居民感知类型及声环境满意度的形成机制，为噪声精准治理提供实证依据。

\textbf{数据与方法：}首先，对问卷进行有效性筛选与质量检验，获得\textbf{467份有效样本}，样本覆盖不同年龄、文化程度、城市及居住区域，信度与效度检验均满足分析要求。其次，采用\textbf{探索性因子分析}将14项噪声影响题项归纳为身心健康与认知干扰、情绪与社会关系影响、典型环境噪声源干扰、生活场景噪声源干扰四个因子，累积方差解释率为\textbf{64.03\%}，并利用\textbf{KMeans++聚类}识别出五类差异化居民群体。进一步，构建\textbf{有序Logistic回归}并计算决策概率与平均边际效应，同时利用\textbf{结构方程模型}检验噪声影响的传导路径。最后，将有序Logistic与随机森林、XGBoost进行对比，检验可解释模型的预测效果。

\textbf{主要结论：}浙江省2020—2024年环境噪声等效声级呈\textbf{“高位波动、2024年反弹”}特征，杭甬温三市中\textbf{杭州居民的不满意概率相对最高}。噪声感知维度中，\textbf{情绪与社会关系影响对满意度的负向作用最强}，住宅配备隔音墙可使“不满意”概率平均下降\textbf{13.76个百分点}。结构方程模型显示，噪声源既直接降低满意度，又经情绪与社会关系渠道间接影响满意度；模型对比中，有序Logistic在准确率、宏平均F1和Kappa指标上优于随机森林与XGBoost。调研还发现，居民\textbf{官方投诉与法律维权使用率偏低}，\textbf{超过六成居民支持强化噪声治理}，诉求集中于法规标准、排放时段、隔音设施和投诉机制。

\textbf{治理建议：}据此，本报告提出差异化治理建议：首先，针对隔音墙对不满意概率的显著抑制作用，优先推进老旧小区、交通干线和敏感场所的隔音改造；其次，围绕交通、施工等主要噪声源，强化源头降噪、夜间管控与智能监测；最后，针对居民官方投诉率偏低的问题，完善投诉反馈和社区协商机制，并重点关注青少年及身心情绪高敏感群体，为浙江省“浙里宁静”专项行动提供实证支撑。

\vskip1ex
\noindent{\zihao{-4}\heiti 关键词：}~城市噪声污染；探索性因子分析；聚类分析；有序Logistic回归；XGBoost；结构方程模型

\end{quotation}

\clearpage

%========================================================================================
%	目录
%========================================================================================
\setcounter{tocdepth}{2}
\tableofcontents
\clearpage
\listoffigures
\clearpage
\listoftables
\clearpage
